% ============================================================================
% LANGENTWURF LE-05b: Reguläre Verben auf -er/-ir
% Spanisch Klasse 7 - Sequenz 5: Mi tiempo libre
% ============================================================================
% Tier: 1 (vollständig)
% Doppelstunde: 3-4 von 8
% Fachfarbe: #c41e3a (Karminrot)
% ============================================================================

\documentclass[11pt,a4paper]{article}

\usepackage[utf8]{inputenc}
\usepackage[T1]{fontenc}
\usepackage[ngerman]{babel}
\usepackage{geometry}
\usepackage{fancyhdr}
\usepackage{lastpage}
\usepackage{xcolor}
\usepackage{tcolorbox}
\usepackage{enumitem}
\usepackage{tabularx}
\usepackage{longtable}
\usepackage{array}
\usepackage{booktabs}
\usepackage{amssymb}

% ============================================================================
% KONFIGURATION
% ============================================================================

\newcommand{\fach}{Spanisch}
\newcommand{\fachfarbe}{c41e3a}
\newcommand{\jahrgangsstufe}{7}
\newcommand{\schulart}{Gesamtschule}
\newcommand{\stundenthema}{Reguläre Verben auf -er/-ir}
\newcommand{\reihenthema}{Mi tiempo libre}
\newcommand{\stundennummer}{05b}
\newcommand{\reihenstunden}{8}
\newcommand{\gerniveau}{A1}
\newcommand{\lehrwerk}{Qué pasa 1}
\newcommand{\lehrwerkseite}{S. 44-46}
\newcommand{\lerngruppengroesse}{24}

% ============================================================================
% LAYOUT
% ============================================================================
\geometry{left=2.5cm, right=2.5cm, top=2.5cm, bottom=2.5cm}
\definecolor{fach}{HTML}{\fachfarbe}
\definecolor{gray}{HTML}{666666}
\definecolor{lightgray}{HTML}{f5f5f5}
\definecolor{successgreen}{HTML}{2a7a4b}
\definecolor{warningorange}{HTML}{b35c00}
\definecolor{basis}{HTML}{2a7a4b}
\definecolor{standard}{HTML}{2c5aa0}
\definecolor{erweiterung}{HTML}{7b1fa2}

\tcbuselibrary{skins,breakable}

\newtcolorbox{infobox}[1][]{
    colback=lightgray,
    colframe=fach,
    fonttitle=\bfseries,
    title=#1,
    breakable
}

\newtcolorbox{scorebox}{
    colback=white,
    colframe=fach,
    fonttitle=\bfseries,
    title=Didaktik-Score,
    breakable
}

\pagestyle{fancy}
\fancyhf{}
\fancyhead[L]{\small\textcolor{gray}{\fach{} -- Jg. \jahrgangsstufe{} -- \stundenthema}}
\fancyhead[R]{\small\textcolor{gray}{LE-\stundennummer{} | Tier 1}}
\fancyfoot[C]{\small\thepage{} / \pageref{LastPage}}
\renewcommand{\headrulewidth}{0.4pt}

% Differenzierungsmarker
\newcommand{\B}{\textcolor{basis}{\textbf{[B]}}}
\newcommand{\Std}{\textcolor{standard}{\textbf{[S]}}}
\newcommand{\E}{\textcolor{erweiterung}{\textbf{[E]}}}

% ============================================================================
% DOKUMENT
% ============================================================================
\begin{document}

% ============================================================================
% DECKBLATT
% ============================================================================
\begin{titlepage}
    \centering
    \vspace*{2cm}
    {\large\textcolor{gray}{\schulart}}\\[2cm]
    {\Huge\textcolor{fach}{\textbf{Langentwurf}}}\\[0.5cm]
    {\large LE-\stundennummer{} | Tier 1 (vollständig)}\\[2cm]
    {\LARGE\textbf{\stundenthema}}\\[0.5cm]
    {\large Sequenz: \reihenthema}\\[2cm]
    \begin{tabular}{rl}
        \textbf{Fach:} & \fach{} (\gerniveau) \\[0.3cm]
        \textbf{Klasse:} & \jahrgangsstufe \\[0.3cm]
        \textbf{Lehrwerk:} & \lehrwerk, \lehrwerkseite \\[0.3cm]
        \textbf{Doppelstunde:} & \stundennummer{} (Std. 3-4 von 8) \\[0.3cm]
    \end{tabular}
    \vfill
\end{titlepage}

\tableofcontents
\newpage

% ============================================================================
% 1. BEDINGUNGSANALYSE
% ============================================================================
\section{Bedingungsanalyse}

\subsection{Lerngruppe}
\begin{tabular}{ll}
    Kursgröße: & \lerngruppengroesse{} SuS \\
    Jahrgangsstufe: & \jahrgangsstufe \\
    Sprachniveau: & \gerniveau{} (GER) -- Anfänger \\
    Fremdsprache: & 2. Fremdsprache (nach Englisch) \\
\end{tabular}

\subsection{Vorwissen aus Sequenz 05a}
\begin{itemize}
    \item Freizeitaktivitäten benennen: el fútbol, la música, el cine...
    \item Vorlieben ausdrücken: Me gusta / No me gusta
    \item Reguläre Verben auf -ar (hablar, escuchar, tocar, bailar, nadar)
    \item Konjugationsendungen -ar: -o, -as, -a, -amos, -áis, -an
\end{itemize}

\subsection{Differenzierung (intern)}
\begin{tabular}{|l|p{3.5cm}|p{3.5cm}|p{3.5cm}|}
\hline
& \B{} \textbf{Basis} & \Std{} \textbf{Standard} & \E{} \textbf{Erweiterung} \\
\hline
\textbf{Ziel} & BR: Rezeption & MR: Produktion mit Hilfe & GY: Freie Produktion \\
\hline
\textbf{-er/-ir} & Tabelle mit Vorgabe & Konjugation selbst & + Sätze bilden \\
\hline
\textbf{Anwendung} & Lückentext & Sätze ergänzen & Eigene Texte schreiben \\
\hline
\end{tabular}

\vspace{0.5em}
\textbf{WICHTIG:} Diese Marker sind NUR für die Lehrkraft!

% ============================================================================
% 2. SACHANALYSE
% ============================================================================
\section{Sachanalyse}

\subsection{Sprachliche Strukturen}

\textbf{Reguläre Verben auf -er (z.B. comer, beber, leer, correr):}
\begin{center}
\begin{tabular}{|l|l|l|}
\hline
\textbf{Person} & \textbf{Pronomen} & \textbf{-er Endung} \\
\hline
1. Sg. & yo & -o \\
2. Sg. & tú & -es \\
3. Sg. & él/ella/usted & -e \\
1. Pl. & nosotros/-as & -emos \\
2. Pl. & vosotros/-as & -éis \\
3. Pl. & ellos/ellas/ustedes & -en \\
\hline
\end{tabular}
\end{center}

\textbf{Reguläre Verben auf -ir (z.B. vivir, escribir, abrir):}
\begin{center}
\begin{tabular}{|l|l|l|}
\hline
\textbf{Person} & \textbf{Pronomen} & \textbf{-ir Endung} \\
\hline
1. Sg. & yo & -o \\
2. Sg. & tú & -es \\
3. Sg. & él/ella/usted & -e \\
1. Pl. & nosotros/-as & -imos \\
2. Pl. & vosotros/-as & -ís \\
3. Pl. & ellos/ellas/ustedes & -en \\
\hline
\end{tabular}
\end{center}

\textbf{Wichtige Vokabeln für diese Stunde:}
\begin{itemize}
    \item \textbf{-er Verben:} comer (essen), beber (trinken), leer (lesen), correr (laufen/rennen)
    \item \textbf{-ir Verben:} vivir (leben/wohnen), escribir (schreiben), abrir (öffnen)
\end{itemize}

\subsection{Vergleich -ar, -er, -ir}

\begin{center}
\begin{tabular}{|l|l|l|l|}
\hline
\textbf{Person} & \textbf{-ar} & \textbf{-er} & \textbf{-ir} \\
\hline
yo & -o & -o & -o \\
tú & -as & -es & -es \\
él/ella & -a & -e & -e \\
nosotros & -amos & -emos & -imos \\
vosotros & -áis & -éis & -ís \\
ellos & -an & -en & -en \\
\hline
\end{tabular}
\end{center}

\textbf{Merke:} Die Endungen von -er und -ir unterscheiden sich nur in der 1. und 2. Person Plural!

\subsection{Typische Interferenzen}
\begin{itemize}
    \item Verwechslung der Endungen -ar/-er: *``Yo como'' (richtig) vs. *``Yo coma'' (falsch)
    \item Falsches Stammwort: *``viver'' statt ``vivir''
    \item Übertragung englischer Strukturen: *``Yo bebo agua'' ohne ``un poco de''
\end{itemize}

% ============================================================================
% 3. DIDAKTISCHE ANALYSE
% ============================================================================
\section{Didaktische Analyse}

\subsection{Stellung in der Sequenz}
\begin{tabular}{|c|l|l|}
\hline
\textbf{Block} & \textbf{Thema} & \textbf{Status} \\
\hline
5a & Freizeit + Verben auf -ar & abgeschlossen \\
\textbf{5b} & \textbf{Verben auf -er/-ir} & \textbf{diese Stunde} \\
5c & Tagesablauf + Uhrzeiten & folgt \\
5d & Mi día perfecto + Sequenztest & folgt \\
\hline
\end{tabular}

\subsection{Kompetenzziele (Rahmenplan MV)}
\begin{itemize}
    \item \textbf{Kommunikativ:} Über Freizeitaktivitäten sprechen, Alltagshandlungen beschreiben
    \item \textbf{Sprachlich:} Grammatik (Verbkonjugation -er/-ir), Wortschatz (Freizeit)
    \item \textbf{Methodisch:} Regelableitung, Analogieschlüsse (Transfer von -ar)
\end{itemize}

\subsection{Legitimation}
\textbf{Rahmenplan Spanisch MV (Kl. 7-10):}
\begin{itemize}
    \item Kompetenzbereich: Kommunikative Kompetenz -- Sprechen und Schreiben
    \item Standard: ``Die SuS können einfache Alltagssituationen sprachlich bewältigen und über Gewohnheiten berichten.''
\end{itemize}

% ============================================================================
% 4. STUNDENZIELE
% ============================================================================
\section{Stundenziele}

\subsection{Grobziel}
Die SuS erweitern ihre \textbf{grammatische Kompetenz}, indem sie reguläre Verben auf -er und -ir konjugieren und in Sätzen über Freizeitaktivitäten anwenden.

\subsection{Feinziele (operationalisiert)}
\begin{tabular}{|c|p{7.5cm}|c|c|}
\hline
\textbf{Nr.} & \textbf{Die SuS können...} & \textbf{AFB} & \textbf{Diff.} \\
\hline
FZ1 & die Konjugationsendungen von -er und -ir Verben \textbf{nennen} & I & alle \\
\hline
FZ2 & regelmäßige -er/-ir Verben in allen Personen korrekt \textbf{konjugieren} & II & alle \\
\hline
FZ3 & den Unterschied zwischen -ar, -er und -ir Endungen \textbf{erklären} & II & \Std{}/\E{} \\
\hline
FZ4 & eigene Sätze über Freizeitaktivitäten mit -er/-ir Verben \textbf{formulieren} & III & \E{} \\
\hline
\end{tabular}

\vspace{1em}
\textbf{AFB-Verteilung:} AFB I: 25\% | AFB II: 50\% | AFB III: 25\%

% ============================================================================
% 5. METHODISCHE ANALYSE
% ============================================================================
\section{Methodische Analyse}

\subsection{Einstieg}
\begin{itemize}
    \item \textbf{Methode:} Bildimpuls: Jugendliche bei Freizeitaktivitäten (lesen, essen, Sport)
    \item \textbf{Begründung:} Aktivierung Vorwissen, Bezug zum Thema ``tiempo libre''
    \item \textbf{Alternative:} Kurzes Video über Tagesablauf
\end{itemize}

\subsection{Erarbeitung}
\begin{itemize}
    \item \textbf{Methode:} Entdeckendes Lernen -- Endungen aus Beispielsätzen ableiten
    \item \textbf{Transfer:} Vergleich mit bereits bekannten -ar Endungen
    \item \textbf{Scaffolding:} Farbcodierung (-er = blau, -ir = grün), Tabelle schrittweise
\end{itemize}

\subsection{Übungsphasen}
\begin{itemize}
    \item \textbf{Methode:} Chorisches Sprechen $\rightarrow$ Partnerarbeit $\rightarrow$ Schreibübung
    \item \textbf{Progression:} Gelenkt $\rightarrow$ Halbfrei $\rightarrow$ Frei
    \item \textbf{Differenzierung:} Hilfestellungen mit/ohne Endungstabelle
\end{itemize}

\subsection{Medien}
\begin{itemize}
    \item Tafel/Whiteboard: Konjugationstabelle
    \item Lehrwerk: \lehrwerk, \lehrwerkseite
    \item Arbeitsblatt: AB-05b.pdf
    \item Lernpfad: LP-05b.html (Tablets, optional)
\end{itemize}

% ============================================================================
% 6. VERLAUFSPLANUNG
% ============================================================================
\section{Verlaufsplanung}

\textbf{1. Stunde (45 Min.)}

\begin{longtable}{|p{1cm}|p{1.5cm}|p{4cm}|p{3cm}|p{1.5cm}|p{2cm}|}
\hline
\textbf{Zeit} & \textbf{Phase} & \textbf{Lehrerhandeln} & \textbf{Schülerhandeln} & \textbf{SF} & \textbf{Material} \\
\hline
\endhead

0--5 & Einstieg & Bilder zeigen: ``¿Qué hacen?'' & Beschreiben mit Vorwissen & Plenum & Bilder \\
\hline

5--10 & Aktivier. & Wiederholung -ar Endungen: ``hablar, escuchar...'' & Konjugation nennen & Plenum & Tafel \\
\hline

10--20 & Input & Verben auf -er einführen: comer, beber, leer, correr. Beispielsätze, Endungen herleiten & Muster erkennen, Tabelle ergänzen & Plenum & Tafel, AB \\
\hline

20--30 & Übung 1 & Chorisches Sprechen: ``Yo como, tú comes...'' & Nachsprechen, Aussprache üben & Plenum & -- \\
\hline

30--40 & Übung 2 & EA: Lückentext -er Verben & AB bearbeiten & EA & AB-05b \\
\hline

40--45 & Sicherung & Lösungen besprechen, Merkkasten -er & Kontrollieren, ergänzen & Plenum & AB-05b \\
\hline
\end{longtable}

\vspace{1em}

\textbf{2. Stunde (45 Min.)}

\begin{longtable}{|p{1cm}|p{1.5cm}|p{4cm}|p{3cm}|p{1.5cm}|p{2cm}|}
\hline
\textbf{Zeit} & \textbf{Phase} & \textbf{Lehrerhandeln} & \textbf{Schülerhandeln} & \textbf{SF} & \textbf{Material} \\
\hline
\endhead

0--5 & Aktivier. & Quiz: -er Endungen abfragen & Antworten & Plenum & -- \\
\hline

5--15 & Input & Verben auf -ir einführen: vivir, escribir, abrir. Vergleich mit -er & Unterschiede erkennen & Plenum & Tafel \\
\hline

15--25 & Übung 3 & PA: Sätze bilden mit -er/-ir Verben & Sätze formulieren, vorlesen & PA & Karten \\
\hline

25--35 & Vertiefung & Vergleichstabelle -ar/-er/-ir. \E{}: Eigener Text ``Mi día'' & Tabelle ausfüllen, schreiben & EA/PA & AB-05b \\
\hline

35--40 & Sicherung & Merkkasten -ir ausfüllen, Unterschiede zusammenfassen & AB ergänzen & Plenum & AB-05b \\
\hline

40--45 & Abschluss & Zusammenfassung, HA: Vokabeln, LP-05b optional & Notieren & Plenum & -- \\
\hline
\end{longtable}

\textbf{Legende:} EA = Einzelarbeit, PA = Partnerarbeit

% ============================================================================
% 7. ERWARTETE SCHWIERIGKEITEN
% ============================================================================
\section{Erwartete Schwierigkeiten}

\begin{tabular}{|p{4.5cm}|p{8cm}|}
\hline
\textbf{Schwierigkeit} & \textbf{Reaktion} \\
\hline
Verwechslung -er/-ir Endungen in nosotros/vosotros & Farbcodierung: -emos/-éis (blau) vs. -imos/-ís (grün) \\
\hline
Transfer von -ar Endungen & Expliziter Vergleich: -as $\rightarrow$ -es, -a $\rightarrow$ -e \\
\hline
Aussprache von ``leer'' [le.'er] & Vorsprechen, Betonung auf zweiter Silbe \\
\hline
Stammwechsel erwartet (gibt es hier nicht) & Beruhigen: Diese Verben sind alle regelmäßig! \\
\hline
Zeitproblem & Puffer: Vergleichstabelle als HA \\
\hline
\end{tabular}

% ============================================================================
% 8. HAUSAUFGABE
% ============================================================================
\section{Hausaufgabe}

\begin{itemize}
    \item Lehrwerk S. 45, Übung 4 (Verben konjugieren)
    \item Vokabeln lernen: comer, beber, leer, correr, vivir, escribir, abrir
    \item Optional: Lernpfad LP-05b.html bearbeiten
\end{itemize}

% ============================================================================
% 9. LÖSUNGEN
% ============================================================================
\section{Lösungen}

\subsection{AB-05b Lösungen}

\textbf{Merkkasten -er:}
\begin{itemize}
    \item yo como, tú comes, él/ella come
    \item nosotros comemos, vosotros coméis, ellos comen
\end{itemize}

\textbf{Merkkasten -ir:}
\begin{itemize}
    \item yo vivo, tú vives, él/ella vive
    \item nosotros vivimos, vosotros vivís, ellos viven
\end{itemize}

\textbf{Aufgabe 1: Konjugation -er}
\begin{enumerate}[label=\alph*)]
    \item Yo \textbf{como} pizza. (comer)
    \item Tú \textbf{bebes} agua. (beber)
    \item María \textbf{lee} un libro. (leer)
    \item Nosotros \textbf{corremos} en el parque. (correr)
\end{enumerate}

\textbf{Aufgabe 2: Konjugation -ir}
\begin{enumerate}[label=\alph*)]
    \item Yo \textbf{vivo} en Rostock. (vivir)
    \item Tú \textbf{escribes} un mensaje. (escribir)
    \item Pedro \textbf{abre} la ventana. (abrir)
    \item Nosotros \textbf{vivimos} en Alemania. (vivir)
\end{enumerate}

\textbf{Aufgabe 3: Sätze bilden}
\begin{itemize}
    \item Beispiel: Yo leo libros en mi tiempo libre.
    \item Beispiel: Mi hermano corre en el parque.
    \item Beispiel: Nosotros comemos pizza los viernes.
\end{itemize}

\textbf{Aufgabe 4: Vergleich (Bewertungskriterien)}
\begin{itemize}
    \item Unterschied nosotros korrekt erkannt: 2P
    \item Unterschied vosotros korrekt erkannt: 2P
    \item Eigene Beispiele korrekt: 2P
\end{itemize}

% ============================================================================
% 10. ANLAGEN
% ============================================================================
\section{Anlagen}

\begin{enumerate}
    \item Arbeitsblatt (AB-05b.pdf)
    \item Musterlösung (ML-05b.pdf)
    \item Lehrerhinweise (LH-05b.pdf)
    \item Lernpfad (LP-05b.html)
    \item Verbkarten für Partnerarbeit (Kopiervorlage)
\end{enumerate}

% ============================================================================
% 11. DIDAKTIK-SCORE
% ============================================================================
\newpage
\section{Didaktik-Score}

\begin{scorebox}
\textbf{Bewertung der didaktischen Qualität nach 10 Dimensionen}\\
\textbf{Ziel-Score: $\geq$ 9.0 (Premium-Qualität)}

\vspace{1em}
\begin{tabular}{|c|l|c|c|p{4cm}|}
\hline
\textbf{\#} & \textbf{Dimension} & \textbf{Gew.} & \textbf{Score} & \textbf{Notiz} \\
\hline
1 & Differenzierung & 15\% & 9/10 & [B]/[S]/[E] qualitativ \\
\hline
2 & Taxonomie (AFB) & 10\% & 9/10 & 25/50/25 erfüllt \\
\hline
3 & Lernziel-Alignment & 10\% & 10/10 & RP MV explizit \\
\hline
4 & Aktivierung & 10\% & 9/10 & Entdeckend, Transfer \\
\hline
5 & Scaffolding & 10\% & 10/10 & Farbcodierung, gestuft \\
\hline
6 & Feedback-Qualität & 10\% & 9/10 & LP: elaboriert \\
\hline
7 & Sprachliche Klarheit & 10\% & 10/10 & Kl. 7 angemessen \\
\hline
8 & Motivation & 10\% & 9/10 & Alltagsbezug Freizeit \\
\hline
9 & Inklusion (UDL) & 10\% & 9/10 & Mehrere Zugänge \\
\hline
10 & Praktikabilität & 5\% & 9/10 & 90 Min. realistisch \\
\hline
\multicolumn{3}{|r|}{\textbf{GESAMT:}} & \textbf{9.3/10} & \\
\hline
\end{tabular}

\vspace{1em}
$\checkmark$ \textcolor{successgreen}{\textbf{FREIGABE} (Score $\geq$ 9.0)}
\end{scorebox}

\end{document}
